\chapter*{Note to Students}
\addcontentsline{toc}{chapter}{Note to Students}

As the title of the book indicates, this is not a book that is just to be read.  It was written
so that the reader interacts with the material.  If you attempt to just read what is 
written and take no part in the exercises that are embedded throughout, you will likely
get very little out of it.  Learning needs to be active, not passive.  The more active you are
as you `read' the book, the more you will get out of it.  That will translate to better learning.
And it will also translate to a higher grade.  So whether you are motivated by learning (which is
my hope) or merely by getting a certain grade, your path will be the same--use this book as 
described below.

The content is presented in the following manner.  
First, concepts and definitions are given--generally one at a time.  
Then one or more examples that illustrate the concept/definition will be given.  
After that you will find one or more exercises of various kinds.  This is where this book
differs from most.  Instead of piling on more examples that you merely read and {\em think} you understand,
you will be asked to solve some for yourself so that you can be more confident that you really {\em do}
understand. 

Some of the exercises are just called {\em Exercises}.  
They are very similar to the examples, except that you have to provide the solution.  
There are also {\em Fill in the details} which provide part of the solution, but ask you to provide
some of the details.  The point of these is to help you think about some of the finer details that
you might otherwise miss.  There are also {\em Questions} of various kinds that get you thinking about
the concepts. Finally, there are {\em Evaluate} exercises.  These ask you to look at solutions written
by others and determine whether or not they are correct.  More precisely, your goal is to try to find
as many errors in the solutions as you can.  Usually there will be one or more errors in each solution,
but occasionally a correct solution will be given, so pay careful attention to every detail.  The point 
of these exercises is to help you see mistakes before you make them.  Many of these exercises are based
on solutions from previous students, so they often represent the common mistakes students make.
Hopefully if you see someone else make these mistakes, you will be less likely to make them yourself. 

The point of the exercises is to get you thinking about and interacting with the material.  
As you encounter these, you should write your solution in the space provided.  After you have written
your solution, you should check your answer with the solution provided in the back of the book.  
You will get the most out of them if you first do your
best to give a complete solution on your own, and then {\em always} check your solution with the one provided
to make sure you did it correctly.  If yours is significantly different, make sure you determine whether or
not the differences are just a matter of choice or if there is something wrong with your solution.

If you get stuck on an exercise, you should re-read the previous material 
(definitions, examples, etc.) and see if that helps.  Then give it a little more thought.  
For {\em Fill in the details} questions, sometimes reading what is past a blank will help you
figure out what to put there.
If you get {\em really} stuck on an exercise, look up the solution and make sure you fully understand it.  
But don't jump to the solution too quickly or too often without giving an honest attempt at solving the exercise
yourself.  When you do end up looking up a solution, you should always try to rewrite it in the
space provided {\em in your own words}.  
You should not just copy it word for word. You won't learn as 
much if you do that.  Instead, do your best to fully understand the solution.  Then, without looking
at the solution, try to re-solve the problem and write your solution in the space provided.  
Then check the solution again to make sure you got it right.

It is highly recommended that you act as your own grader when you check your solutions.
If your solution is correct, put a big check mark in the margin.  
If there are just a few errors, 
use a different colored writing utensil to mark and fix your errors.
If your solution is way off, cross it out (just put a big `X' through it) 
and write out your second attempt, using a separate sheet of paper if necessary. 
If you couldn't get very far without reading the solution, you should somehow indicate that. 
So that you can track your errors, I highly recommend crossing out incorrect solutions
(or portions of solutions) instead of erasing them.  
Doing this will also allow you to look back and determine how well you did as you were working
through each chapter.  
It may also help you determine how to spend your time as you study for exams. 

This whole process will help you become better at evaluating your own work. 
This is important because you should be confident in your answers, but only when they
are correct.  Grading yourself will help you gain confidence when you are correct and
help you quickly realize when you are not correct so that you do not
become confident about the wrong things.
Another reason that grading your solutions is important is so that when you go back
to re-read any portion of the book, you will know whether or not what you wrote
was correct.

It is important that you read the solutions to the exercises after you attempt them, even if you
think your solution is correct.  
The solutions often provide further insight into the material and should be regarded 
as part of any reading assignment given.

Make sure you read carefully.  When you come upon an {\em Evaluate} exercise, do not
mistake it for an example.  Doing so might lead you down the wrong path.  Never consider
the content of an {\em Evaluate} exercise to be correct unless you have verified with the 
solution that it is really correct.  To be safe, when re-reading, always assume that the
{\em Evaluate} exercises are incorrect, and never use them as a model for your own problem 
solving.  To help you, we have tried to differentiate these from other example and exercise
types by using a different font.

There is an expectation that you are able to solve every exercise on your own.
(Note that I am talking about the exercises embedded into the chapters, 
not the homework problems at the end of each chapter.)
If there are exercises that you are unable to complete, you need to get them cleared up
immediately.  This might mean asking about them in class, going to see the professor or a
teaching assistant, and/or going to a help center/tutor.  Whatever it takes, make sure you 
have a clear understanding of how to solve all of them.  

Every chapter ends with two sections called {\em Reading Comprehension Questions}
and {\em Problems}. The {\em Problems} sections are exactly what they sound like--a list of problems
suitable for working on in class or given as homework assignments.

All of the {\em Reading Comprehension Questions} should be attempted after
you have finished reading each section (including completing all of the exercises). 
They are sort of the final check of your comprehension of the material before you move on to solving 
homework problems.
Although some of these questions are similar to the exercises in the sections,
others are more conceptual in nature. The majority of them are not meant to be difficult, but rather to
test whether you really understand the material from the section as whole. These can be used as a starting point
for class discussion, so be sure to ask about those that you have trouble completing and/or are unsure about.

Space is not given in the book for solutions to the {\em Reading Comprehension Questions}, 
so write your answers on paper or use a Google Doc or other typesetting software to record your solutions.
(In my classes I have students share a Google Doc with me in which they place their answers to these questions,
adding the most recent answers at the top of the document to make it easier to find their recent answers.
When questions can't easily be done in a Google Doc, they write their solution on paper, scan or take a picture of it, 
and include the picture in their Google Doc.)

Solutions to the {\em Reading Comprehension Questions} are available in the back of the book.
As with the exercises throughout the book, it is highly recommended that you check your answers and grade 
your own work, crossing out your solution when you were incorrect (instead of erasing/deleting it)
and replacing it with the correct solution.


